\chapter{What the existing \textit{Senecio} results already contribute}

\section{Local adaptation and early reproductive isolation}

Reciprocal transplants across Dune, Headland, Tableland, and Woodland environments showed strong local adaptation and substantial immigrant inviability. F1 hybrids exhibited heterosis, yet hybrids with a native parent performed better than hybrids with two foreign parents. F2 hybrids showed a pronounced fitness reduction in a common garden, consistent with early intrinsic incompatibilities emerging after recombination \cite{walter2016}.

\begin{biobox}
This combination is exactly the biological setting in which an early geometric mechanism is plausible: strong environment-dependent barriers are already present, F1 vigor can mask incompatibility, and recombinant generations reveal additional failure.
\end{biobox}

\section{Evolution of genetic variance}

The covariance-tensor analysis of ecotype \(\G\) matrices found divergence in the magnitude and multivariate distribution of genetic variance. Divergence in \(\G\), including divergence in the orientation of \(\mathbf g_{\max}\), aligned with divergence in mean phenotype \cite{walter2018}.

\begin{connectionbox}
That result supplies empirical precedent for the first half of the geometric snowball: divergent natural selection can be associated with changes in the distribution and orientation of standing genetic variance. The new theory adds local fitness curvature and asks whether the resulting \(\Sop\) subspaces predict hybrid load.
\end{connectionbox}

\section{Loss of ecologically important variation after hybrid breakdown}

In later-generation hybrids, severe F2 fertility loss was followed by fitness recovery, reduced variance in field performance, loss of parental fitness trade-offs, and loss of phenotypic extremes in the F4. These patterns suggest that variation associated with adaptation was also involved in negative interactions contributing to intrinsic isolation \cite{walter2020}.

\begin{connectionbox}
The trace identity gives a continuous quantitative version of this idea. It asks whether covariance adapted to one curvature matrix becomes costly when averaged or recombined with covariance adapted to another. The DMI layer then captures additional locus-specific negative interactions that need not be visible in additive covariance alone.
\end{connectionbox}

\section{The replicated Dune--Headland opportunity}

The strongest test is pair-specific rather than ecotype-pooled. Replicate Dune--Headland pairs allow environmental repeatability to be separated from pair-specific mutation order. The theory predicts that repeatable habitat forcing should contribute a repeatable component to \(\dot\Sop\) and \(\cK\), whereas stochastic histories should generate pair-specific departures.

\begin{figure}[H]
\centering
\includegraphics[width=0.90\textwidth]{figures/senecio_pair_descriptive.pdf}
\caption{Exploratory pair-level summary from available morphology and soil data. Both axes are descriptive distances only. The weak association is informative because it reinforces that soil-chemistry distance cannot substitute for directly estimated fitness curvature and morphology-mean distance cannot substitute for channel angle.}
\label{fig:senecioExploratory}
\end{figure}
